\documentclass{article}
\begin{document}
A typesetting system takes a description of a document and decides where every character goes. It breaks paragraphs into lines and stacks the lines into pages. Lines of a paragraph are chosen together rather than one at a time, so that the spacing is as even as the words allow and few hyphens appear. A typesetting system takes a description of a document and decides where every character goes. It breaks paragraphs into lines and stacks the lines into pages.

\begin{figure}[ht]
\rule{4cm}{2cm}
\caption{A black box standing for a figure.}
\end{figure}

Lines of a paragraph are chosen together rather than one at a time, so that the spacing is as even as the words allow and few hyphens appear. A typesetting system takes a description of a document and decides where every character goes. It breaks paragraphs into lines and stacks the lines into pages. Lines of a paragraph are chosen together rather than one at a time, so that the spacing is as even as the words allow and few hyphens appear.

\begin{table}[t]
\centering
\begin{tabular}{lr}
a & 1 \\
b & 22 \\
\end{tabular}
\caption{A small table.}
\end{table}

A typesetting system takes a description of a document and decides where every character goes. It breaks paragraphs into lines and stacks the lines into pages. Lines of a paragraph are chosen together rather than one at a time, so that the spacing is as even as the words allow and few hyphens appear. A typesetting system takes a description of a document and decides where every character goes. It breaks paragraphs into lines and stacks the lines into pages. Lines of a paragraph are chosen together rather than one at a time, so that the spacing is as even as the words allow and few hyphens appear.
\end{document}
